\section{Status and limitations}
\label{sec:status}

\subsection{Testing and continuous integration}

The unit tests run without a language model. The ingestion tests check the
PDF, Word and text readers and the chunk sizes. The generation tests replace
Ollama with a stub that records the prompt, and check that the answer and
detail are taken from the model's JSON reply, that a reply which is not JSON
is kept as the answer, that token usage is reported, that the focus line
appears exactly when a named provision was found, and that interleaving keeps
the strongest passage first. A pipeline test ingests a PDF fixture and checks
that retrieval returns passages with their locations; it needs the embedding
model and the vector store and is run locally.

Continuous integration runs on GitHub Actions for every pull request and every
push to the main branch. One job installs the interface's dependencies,
type-checks it and builds it; the other installs the backend on Python 3.12
with a CPU-only build of PyTorch and runs the eleven model-free tests.

\subsection{Measured behaviour}

\Cref{tab:corpora} lists what ingestion recovers from the two main development
documents, and \cref{tab:measured} how representative questions are served.
Times were measured on the development machine described in
\cref{sec:overview}.

\begin{table}[ht]
\centering
\small
\begin{tabularx}{\textwidth}{@{}lrrX@{}}
\toprule
Document & Pages & Chunks & Structure found \\
\midrule
AI Act~\cite{euaiact} & 144 & 2\,019 & 13 chapters, 5 sections, 113 articles,
                                       180 recitals, 13 annexes; printed pages
                                       equal PDF pages \\
TCPS~2 (2022)~\cite{tcps2} & 287 & 1\,664 & 13 chapters, 142 articles; printed
                                       page = PDF page $-$ 8 \\
\bottomrule
\end{tabularx}
\caption{Structure recovered at ingestion.}
\label{tab:corpora}
\end{table}

\begin{table}[ht]
\centering
\small
\begin{tabularx}{\textwidth}{@{}Xlrr@{}}
\toprule
Question & Path & Passages & Retrieval \\
\midrule
``paragraph 2 article 7'' (AI Act) & exact, subdivision & 14 & 3.1\,ms \\
``chapter 3 page 48'' (TCPS~2) & exact, chapter and page & 7 & 5.5\,ms \\
``paragraph 2 article 85'' (AI Act) & exact, part missing & 2 & -- \\
``who can lodge a complaint with a market surveillance authority'' (AI Act)
  & hybrid & 5 & 235\,ms \\
\bottomrule
\end{tabularx}
\caption{Retrieval for representative questions. For the missing paragraph
the reply is composed without calling the model; for the hybrid question the
provision on the right to lodge a complaint, Article 85, is ranked first.
Generating an answer adds 7--9\,s once the model is loaded.}
\label{tab:measured}
\end{table}

\subsection{Limitations}

\begin{itemize}
  \item Structure is recognised from the numbering and punctuation
  conventions of formal drafting. A document that does not follow them is
  chunked without labels and searched by similarity only.
  \item A question is resolved to one provision at a time. A question relating
  two provisions (``what does paragraph 3 of this article have to do with
  Article 7'') retrieves only one of them.
  \item References to other instruments are recognised and marked, but not
  followed into other uploaded documents.
  \item The local 3B model and its 4\,096-token context window limit the
  passages to about 2\,000 words and keep summaries short; an amount or a
  condition may be left out of a plain-language summary.
  \item Printed page numbers are recovered only from PDF files with one
  sequential numbering; Word and text files have no pages, and scanned PDF
  files without a text layer are not read.
  \item The accuracy of answers has not yet been quantified on an annotated
  set of questions.
\end{itemize}

\subsection{Future work}

\begin{itemize}
  \item Retrieve every provision a question names, resolve references such as
  ``this article'' from the conversation, and give the model the
  cross-reference path between the provisions, so that relational questions
  can be answered.
  \item Build an evaluation set of annotated questions and measure retrieval
  (whether the expected provision is among the passages) and faithfulness
  (whether every claim in the answer is supported by a cited passage).
  \item Add a verification step that checks each sentence of an answer
  against the passages it cites.
  \item Follow references to other instruments when those instruments are
  among the uploaded documents.
\end{itemize}

\subsection{Summary}

The system implements the RAG technique as a Python pipeline that answers
questions grounded in uploaded documents and cites its sources. Every answer
is drawn from retrieved passages, and every passage carries its document,
printed page and structural reference. What distinguishes it from a generic
pipeline is that it treats document structure as data. A question that names
a location is answered by exact lookup instead of similarity, returning exactly
the part named---or saying that it does not exist. The number of passages
follows from the question instead of a setting. And the cross-references
between provisions are extracted into a graph that the reader can follow from
any answer, through a deliberately simple reference tree.
